\documentclass[12pt,a4paper]{article}
\usepackage[utf8]{inputenc}
\usepackage[T1]{fontenc}
\usepackage[english]{babel}
\usepackage{amsmath,amssymb,amsthm}
\usepackage{geometry}
\usepackage{hyperref}
\geometry{margin=2.5cm}

\newtheorem{theorem}{Theorem}
\newtheorem{proposition}{Proposition}
\newtheorem{lemma}{Lemma}
\newtheorem{corollary}{Corollary}
\newtheorem{axiom}{Axiom}
\newtheorem{definition}{Definition}

\title{%
  \textbf{Chapter 2 — English Version}\\[0.5em]
  \Large Harmonic Mechanics of Discrete Chaos\\[0.5em]
  \large Mathematical Theory: The Universe is Squared\\[0.5em]
  \normalsize Philippe Thomas Savard, 2026
}
\author{Philippe Thomas Savard}
\date{2026}

\begin{document}
\maketitle
\tableofcontents
\newpage

\begin{abstract}
This chapter develops the harmonic mechanics of discrete chaos, the second pillar
of Savard's theory. We show that the apparent irregularities in the distribution
of prime numbers obey an underlying harmonic mechanics that transforms apparent
chaos into deep geometric order. The theory reveals that the ``discrete chaos''
of primes is in reality a superposition of harmonics whose frequencies are determined
by the spectral geometry of Chapter~1.
\end{abstract}

\section{Introduction: From Chaos to Order}

At first glance, prime numbers seem to be distributed chaotically among the integers.
Yet, as Savard reveals, this distribution obeys a precise \emph{harmonic mechanics}.
The ``chaos'' is merely the discrete manifestation of an underlying continuous system.

\begin{definition}[Discrete Chaos]
We call \emph{discrete chaos} the irregularity observed in the sequence of gaps
between consecutive prime numbers:
\[
  g_k = P_{k+1} - P_k, \quad k \geq 1.
\]
This sequence $(g_k)_{k \geq 1}$ constitutes the discrete manifestation of harmonic chaos.
\end{definition}

\section{Harmonics of the Prime Distribution}

\subsection{Harmonic Decomposition}

The prime counting function $\pi(x)$ admits a harmonic decomposition via Riemann's
explicit formula:
\[
  \pi(x) = \mathrm{li}(x) - \sum_{\rho} \mathrm{li}(x^\rho) - \ln 2 +
  \int_x^{\infty} \frac{dt}{t(t^2-1)\ln t},
\]
where the sum runs over all non-trivial zeros $\rho$ of $\zeta(s)$.

Each term $\mathrm{li}(x^\rho)$ represents a \emph{harmonic} of the prime distribution.
Savard's harmonic mechanics analyzes how these harmonics interact to produce the
observed discrete chaos.

\begin{theorem}[Harmonic Mechanics]
The harmonics of the prime distribution form a mechanical system whose state is
entirely determined by sequences $A$ and $B$ from Chapter~1. In particular, the
amplitude of the $k$-th harmonic is proportional to $1/k$.
\end{theorem}

\section{Geometric Origin of Discrete Chaos}

\begin{proposition}[Geometric Origin of Chaos]
The gaps $g_k = P_{k+1} - P_k$ between consecutive primes are entirely determined
by the spectral geometry of sequences $A$ and $B$. More precisely:
\[
  g_k \equiv 0 \pmod{2}, \quad k \geq 2,
\]
and the distribution of gaps obeys Cramér's law:
\[
  \Pr[g_k > t \cdot \ln^2 P_k] \approx e^{-t}.
\]
\end{proposition}

\section{Unity of the Theory}

\begin{theorem}[Unity of the Theory]
The harmonics of the prime distribution, defined by the zeros of $\zeta(s)$, correspond
exactly to the geometric oscillations of the spectrum of sequences $A$ and $B$ around
the critical line $\Re(s) = 1/2$.
In other words, the harmonic mechanics of discrete chaos and the geometry of the prime
number spectrum are two equivalent descriptions of the same phenomenon.
\end{theorem}

\section{Conclusion}

The harmonic mechanics of discrete chaos reveals that the apparent irregularity in
the distribution of prime numbers hides a precise harmonic structure, entirely coherent
with the spectral geometry developed in Chapter~1. The discrete ``chaos'' is the visible
manifestation of a deep geometric order, accessible through Savard's theory.

\bigskip
\noindent\textbf{Author:} Philippe Thomas Savard\\
\noindent\textbf{Since:} 2016\\
\noindent\textbf{Repository:} Mathematical Theory — The Universe is Squared, 2026

\end{document}
